% !TITLE 十五从军征
% !AUTHOR 无名氏
% !DYNASTY 两汉
% !GENRE 乐府诗
% !ORDER 5

\begin{poem}
十五从军征，八十始得归。
道逢乡里人："家中有阿谁\textsuperscript{1}？"
"遥看是君家，松柏冢累累\textsuperscript{2}。"
兔从狗窦\textsuperscript{3}入，雉\textsuperscript{4}从梁上飞。
中庭生旅谷，井上生旅葵\textsuperscript{5}。
舂谷持作饭，采葵持作羹。
羹饭一时熟，不知贻\textsuperscript{6}阿谁。
出门东向看，泪落沾我衣。
\end{poem}

\begin{annotations}
\notetext{1}{阿谁：谁。阿，汉代口语常用前缀，无实意。}
\notetext{2}{冢累累：坟墓一座连着一座，非常多。冢（zhǒng），坟墓。累累，累积聚集的样子。}
\notetext{3}{狗窦：狗洞，给狗进出的小门。窦（dòu），孔洞。}
\notetext{4}{雉：雉（zhì），野鸡。}
\notetext{5}{旅谷、旅葵：野生的谷子和野生的冬葵（一种蔬菜）。旅，野生，指不经播种而自生的谷物草木。}
\notetext{6}{贻：贻（yí），赠送，送给。这里指送饭给亲人吃。}
\end{annotations}

\begin{appreciation}
这首诗是汉乐府中最沉痛的篇章之一，写的是一个老兵六十五年军旅生涯后的归家惨况。

"十五从军征，八十始得归"——十五岁参军，八十岁才回家。六十五年的跨度，一个人从少年变成了老人。这第一句话就给出了一个令人窒息的时间尺度。汉代征兵制度，理论上服役到一定年限可以还乡，但战争连年不断，许多人被长期滞留边疆，一去就是一生。

"道逢乡里人：'家中有阿谁？'"——路上遇到同乡，急切地问："家里还有谁？""阿谁"是汉代的口语，就是"谁"。这个问句暴露了他最深的恐惧：六十五年过去了，家人还在吗？

"遥看是君家，松柏冢累累"——远远看去那就是你家，只是松树柏树下面全是坟墓。同乡没有直接回答，而是让他自己看。这个回答残酷到了极致：家还在，但人已经不在了。"冢累累"——坟墓一座挨着一座，暗示全家都已去世。这种间接的告知，比直接说"你家人都死了"更令人心碎。

接下来的六句，写老兵回到家中看到的景象。"兔从狗窦入，雉从梁上飞"——野兔从狗洞里钻进钻出，野鸡在屋梁上飞来飞去。家已经荒废成了野兽的巢穴。"中庭生旅谷，井上生旅葵"——院子里长满了野生的谷子，井台上长满了野生的冬葵。"旅"就是野生的意思。人不在了，植物就占据了人的空间。这种荒凉不是刻意的描写，而是自然发生的——自然不在乎人类的悲欢。

"舂谷持作饭，采葵持作羹"——他舂了谷子做饭，采了冬葵做汤。这是本能的行为：回家就要做饭。但"羹饭一时熟，不知贻阿谁"——饭做好了，却不知道送给谁。这个细节是全诗最痛的地方。他六十五年来一直梦想着回家吃一顿热饭，现在饭做好了，却没有人可以分享。

"出门东向看，泪落沾我衣"——他走出门，向东望去，眼泪落满了衣襟。东向而望，也许是望向亲人坟墓的方向，也许是习惯性地望向太阳升起的方向——那曾经是他出征的方向。这个结尾没有控诉，没有嚎啕，只有一个老人的沉默和泪水。这种克制的悲痛，比任何激烈的言辞都更有力量。
\end{appreciation}
